% results.tex
\section{Main results}\label{res}
\begin{figure}[htbp!]\centering
\caption{Anti-competitive behavior around the winning threshold}\label{fig:anti-comp}
\subfloat[Incumbency]{\includegraphics[width=0.8\textwidth]{incumbent.png}}\\
\subfloat[Backlog (standardized, 90-day)]{\includegraphics[width=0.8\textwidth]{logrol_std.png}}\\
\begin{minipage}{0.8\textwidth}
{\scriptsize \emph{Notes:} This figure presents graphical results from the regression discontinuity analysis. The horizontal axis in each panel is the margin of victory (running variable), measured as defined in the text, where negative values correspond to cases where a firm won the auction (by bidding just below the closest competitor) and positive values correspond to cases where the firm lost (its bid was just above the winning bid). Panel (a) plots the average incumbency rate (probability the firm was the winner of its previous auction) for bids around the threshold, and panel (b) plots the average standardized backlog (past 90-day contract total) for those bids. The discontinuities at the cutoff (margin = 0) indicate the estimated jump in each outcome for firms that just won compared to those that just lost. \par}
\end{minipage}
\end{figure}

We begin by examining the evidence for anti-competitive behavior in our sample using the RDD approach described above. Figure~\ref{fig:anti-comp} provides a visual summary of the main results, displaying the relationship between the margin of victory and two key outcomes (incumbency and backlog) near the zero cutoff. In panel (a), we observe a clear downward jump in the incumbency rate at the threshold: firms that narrowly won an auction have a substantially higher probability of having won their prior auction than firms that narrowly lost. In panel (b), a similar discontinuity is evident for the backlog measure: narrowly winning firms tend to have a significantly greater recent backlog of contract awards compared to narrowly losing firms. These graphical patterns suggest a marked advantage for firms on the winning side of the threshold.

To quantify these effects, we estimate RDD models for each outcome. The estimated discontinuity (at $\Delta = 0$) in the incumbency outcome is about $-0.171$ (with a standard error of $0.009$), indicating that a firm that just wins an auction is about 17.1 percentage points more likely to have been the incumbent going into that auction than a similarly positioned firm that just loses. Likewise, for the backlog outcome, we estimate a discontinuity of approximately $-0.271$ (std.~err.~$0.017$), implying that narrow winners have a backlog about 0.27 standard deviation units higher than narrow losers.\footnote{These estimates are obtained using a local linear RDD implemented with the \texttt{rdrobust} package in Stata, employing a bias-corrected estimator and clustering at the auction level to account for the pairing of bids within each auction. The results are robust to using different bandwidth selectors and polynomial orders.} Both effects are precisely estimated and statistically significant at the 1\% level. 

The magnitude and direction of these discontinuities carry important implications. In a purely competitive environment, we would not expect the act of barely winning an auction to be strongly associated with having won before or with holding a larger backlog of ongoing projects. The fact that we do observe such associations suggests that an external factor—likely related to favoritism or some form of collusion—is conferring a consistent advantage to certain firms. Notably, the pattern we uncover aligns more closely with \textit{preferential treatment} or corrupt favoritism rather than classical cartel behavior. In a bid rotation cartel, one would expect losing firms to also have had recent successes (since cartel members take turns winning contracts), which would mean no large gap in backlog or incumbency between narrow winners and losers. Instead, our findings show the opposite: the firms that win tend to keep winning repeatedly, whereas those that lose remain with relatively little recent business. This asymmetry points toward a scenario in which particular firms—potentially those with political connections or insider status—are being systematically favored in the procurement process. Such favoritism allows them to accumulate successive wins and a greater share of contracts over time, at the expense of competitors who do not enjoy the same privilege.
